\section{関連研究}
\label{sec_related}
\subsection{大喜利とユーモアの研究}
大喜利を対象とした研究として，Zhong らは画像やテキストのお題に対する大喜利の回答を集めたデータセット Oogiri-GO を構築し，大規模言語モデルによるユーモアの生成を評価した\cite{ref_oogiri_go}．
Hessel らは，風刺漫画のキャプションのコンテストのデータを用い，ユーモアの理解を，キャプションと漫画の対応付け，順位付け，説明の3つの課題として評価した\cite{ref_newyorker}．
これらの研究は回答の面白さを主な対象としており，回答がどのような作りの文であるかを表すラベルは扱っていない．
本研究は，面白さそのものではなく，回答の構造（制約の順守，発話の型，前提とする知識）を，生成や判定の条件や分析の軸として付与することを目的とする．

\subsection{視覚言語モデルと知識蒸留}
近年の VLM は，画像と指示文を入力として，指示に従った回答を学習なしで生成できる．
本研究では，Qwen3-VL\cite{ref_qwen3vl}，InternVL3.5\cite{ref_internvl35}，Gemma 3\cite{ref_gemma3}，SmolVLM\cite{ref_smolvlm}，LFM2-VL\cite{ref_lfm2vl}の公開モデルを比較した．
大規模なモデルの推論には vLLM\cite{ref_vllm}を用いた．
大規模なモデル（教師）の出力を小さなモデル（生徒）に学習させる知識蒸留\cite{ref_kd}は，推論の費用を下げる一般的な方法である．
本研究では，教師ラベルで学習した後に少量の目視ラベルで微調整する二段階の学習を用いる．
これは，教師ラベルの誤りを目視ラベルで補正する点で，ラベルの雑音を含む学習\cite{ref_noisy_student}と関係する．
大規模なモデルの微調整には，低ランクの行列を追加して学習する LoRA\cite{ref_lora}を用いた．
文の埋め込みには，日本語の文埋め込みモデル ruri-v3\cite{ref_ruri}と多言語の multilingual-e5\cite{ref_me5}を用いた．

\subsection{推論の高速化}
言語モデルの推論では，入力の先頭部分（プレフィックス）が共通する複数の要求の間で，注意機構のキーと値（KV キャッシュ）を共有することで計算を省く方法が用いられている\cite{ref_vllm, ref_sglang}．
本研究では，同じお題画像に投稿されたボケが，画像とお題の部分を共通のプレフィックスとして持つことに着目し，分類器の推論にこの考え方を適用する．
